%! Change in Linear and Exponential Functions — Unit 4, Lesson 4.2 — Lesson Plan
\documentclass[10pt]{article}
\usepackage{precalculus-article}
\usepackage{precalculus-boxes}
\usepackage{graphicx}
\graphicspath{{images/}}
\newcommand{\TallMath}[1]{$\displaystyle #1\rule[-1.4em]{0pt}{3.2em}$}
\newcommand{\CourseName}{Precalculus}
\newcommand{\MeetingLength}{55 minutes}
\newcommand{\UnitNumberName}{Unit 4: Exponential Functions \quad}
\newcommand{\LessonNumberName}{Lesson 4.2: Change in Linear and Exponential Functions}

\begin{document}
\begin{center}
    {\Large\bfseries \CourseName \\
    {\small\bfseries \UnitNumberName \LessonNumberName}}
\end{center}

\begin{tcolorbox}[colback=sky, colframe=navy, boxrule=0.9pt, arc=2mm,
    left=2mm, right=2mm, top=1.5mm, bottom=1.5mm]
    \textbf{Primary Objective:} Students will construct linear and exponential functions from a
    point and a rate of change (or ratio), connect these to arithmetic and geometric sequences,
    and determine whether a function is linear or exponential by analyzing output values over
    equal-length input intervals.
    \hfill\textit{\small Skills: 1.C and 3.B}
\end{tcolorbox}

\renewcommand{\arraystretch}{1.6}
\begin{skillbox}[Priority Ideas \& Skills]{goldbox}
\begin{tabularx}{\linewidth}{@{}X X@{}}
\textbf{AP Skills} &
\textbf{Key Understandings} \\
\midrule
\textbf{Skill 1.C} --- Identify a re-expression of mathematical information presented in a given representation. &
\begin{itemize}[leftmargin=1.2em,itemsep=2pt,topsep=0pt]
  \item Linear functions $f(x) = b + mx$ extend arithmetic sequences; exponential functions
        $f(x) = ab^x$ extend geometric sequences. (EK 2.2.A.1, 2.2.A.3)
  \item Both types can be built from one point and a constant rate or ratio.
        (EK 2.2.A.2, 2.2.A.4)
  \item Sequences and functions may differ in domain. (EK 2.2.A.5)
\end{itemize} \\
\textbf{Skill 3.B} --- Apply an appropriate mathematical definition, theorem, or test. &
\begin{itemize}[leftmargin=1.2em,itemsep=2pt,topsep=0pt]
  \item Equal-length interval test: constant difference $\Rightarrow$ linear; constant ratio
        $\Rightarrow$ exponential. (EK 2.2.B.1)
  \item Both function families are fully determined by two distinct input-output pairs.
        (EK 2.2.B.3)
\end{itemize} \\
\end{tabularx}
\end{skillbox}

\begin{skillbox}[{Vocabulary, Concepts, \& Theorems}]{greenbox}
\renewcommand{\arraystretch}{1.4}
\begin{tabularx}{\linewidth}{@{} >{\leavevmode\bfseries\color{navy}}p{0.27\linewidth} X @{}}
Linear function        & $f(x) = b + mx$; output changes by a constant amount (slope $m$) over every unit increase in input. \\
Exponential function   & $f(x) = ab^x$; output is multiplied by a constant base $b$ over every unit increase in input. \\
Constant rate of change & The slope $m$ of a linear function; equal to $\Delta y / \Delta x$ over any interval. \\
Constant ratio         & The base $b$ of an exponential function; equal to $y_{n+1}/y_n$ over equal-length intervals. \\
Initial value          & The output when $x = 0$; equals $b$ for $f(x)=b+mx$ and $a$ for $f(x)=ab^x$. \\
Domain                 & The set of allowable inputs; sequences use $\{0,1,2,\ldots\}$ while the corresponding functions use all reals (or a subset). \\
\end{tabularx}
\end{skillbox}

\begin{skillbox}[Activate Prior Knowledge \& Spiral Review]{sky}
    Spiral review (warm-up) covers skills from Lesson 2.1: identifying the common difference $d$
    and writing $a_n$ for an arithmetic sequence; identifying the common ratio $r$ and writing
    $g_n$ for a geometric sequence; and determining sequence type and initial term from two given
    terms. Students connect these sequence ideas directly to the function constructions in today's
    lesson.
\end{skillbox}

\begin{skillbox}[Hook]{sky}
Present two unlabeled tables on the board (no function type announced):

\medskip
\begin{center}
\renewcommand{\arraystretch}{1.4}
\begin{tabular}{c|c}
$x$ & $f(x)$ \\ \hline
0 & 3 \\ 1 & 7 \\ 2 & 11 \\ 3 & 15
\end{tabular}
\qquad\qquad
\begin{tabular}{c|c}
$x$ & $g(x)$ \\ \hline
0 & 2 \\ 1 & 6 \\ 2 & 18 \\ 3 & 54
\end{tabular}
\end{center}

\medskip
Ask: \textbf{``What is the pattern in Table 1? In Table 2?''} After student discussion, ask:
\textbf{``Can you write a rule that predicts the output for any input?''} Today's goal is to
answer that question — and to know \textit{which} type of function to reach for.
\end{skillbox}

\begin{skillbox}[Lesson]{sky}
\begin{multicols}{2}
\textbf{Part 1 (15 min): Linear Functions from Arithmetic Sequences.}

Arithmetic sequence $a_n = a_0 + dn$ uses an initial value plus repeated addition.
Linear functions do the same: $f(x) = b + mx$.

Point-slope form: if we know slope $m$ and one point $(x_i, y_i)$,
\[
  f(x) = y_i + m(x - x_i).
\]
This is the function analogue of $a_n = a_i + d(n - i)$.

Model: given $m = 4$ and the point $(1, 7)$:
$f(x) = 7 + 4(x - 1) = 4x + 3$. Check: $f(0)=3$, $f(2)=11$.

Key difference: sequences have domain $\{0,1,2,\ldots\}$; the function $f$ is defined for all
real $x$.

\columnbreak

\textbf{Part 2 (15 min): Exponential Functions from Geometric Sequences.}

Geometric sequence $g_n = g_0 r^n$ uses an initial value times repeated multiplication.
Exponential functions do the same: $f(x) = ab^x$.

Point-ratio form: if we know ratio $r$ and one point $(x_i, y_i)$,
\[
  f(x) = y_i \cdot r^{x - x_i}.
\]
Model: given $r = 3$ and the point $(0, 2)$:
$f(x) = 2 \cdot 3^{x-0} = 2 \cdot 3^x$. Check: $f(1)=6$, $f(2)=18$.

\textbf{Part 3 (10 min): Comparing Linear and Exponential.}

Equal-length interval test (EK 2.2.B.1):
\begin{itemize}[leftmargin=1em,itemsep=1pt,topsep=0pt]
  \item Constant first difference $\Rightarrow$ linear
  \item Constant ratio of consecutive outputs $\Rightarrow$ exponential
\end{itemize}

Two points uniquely determine each type (EK 2.2.B.3). With $(x_1,y_1)$ and $(x_2,y_2)$:
linear gives $m = \frac{y_2-y_1}{x_2-x_1}$; exponential gives
$b = \left(\frac{y_2}{y_1}\right)^{1/(x_2-x_1)}$.
\end{multicols}
\end{skillbox}

\begin{skillbox}[Active Monitoring]{redbox}
\begin{itemize}[leftmargin=1.2em,itemsep=2pt,topsep=0pt]
  \item After Part 1: ``If I give you the point $(3, 11)$ and tell you the slope is $2$, what is
        $f(0)$?'' (Students must use point-slope form, not assume $f(0) = 11$.)
  \item During Part 1: watch for students who confuse slope with the $y$-intercept; confirm
        $b = f(0)$ only when you simplify, not at the start.
  \item After Part 2: ``What is the base of the exponential function through $(0,2)$ and $(3,54)$?
        How did you find it?'' (Expect $b = 3$ from $54/2 = 27 = 3^3$.)
  \item During Part 3: cold-call with a table — ``Is this linear, exponential, or neither? Show
        your check.'' Confirm students subtract (not divide) for the first-difference test.
\end{itemize}
\end{skillbox}

\begin{skillbox}[Group Work \& Differentiation]{redbox}
\begin{itemize}[leftmargin=1.2em,itemsep=2pt,topsep=0pt]
  \item \textbf{Tier R (Remediate):} Given a completed table, compute first differences and
        ratios; identify the function type; read off $m$ or $r$ and the initial value; write the
        formula using the provided template.
  \item \textbf{Tier A (Apply):} Given two points, determine whether the relationship is linear or
        exponential and construct the function; apply to a contextual scenario with interpretation.
  \item \textbf{Tier E (Extend):} Investigate a population table with an apparent vertical shift;
        determine whether subtracting a constant reveals an exponential pattern; connect to the
        idea of exponential-plus-constant (preview of EK 2.5.A.2).
\end{itemize}
\end{skillbox}

\begin{skillbox}[Individual Work \& Assessment]{redbox}
Exit ticket (3 questions, 1 page): (1) given a table with inputs 0--3 and outputs doubling,
identify the function type and write the formula; (2) construct an exponential function through
$(0,3)$ and $(2,75)$; (3) true/false --- two input-output pairs uniquely determine both a
linear and an exponential function.

Collect and sort: students missing item 1 need interval-test re-teaching (LO 2.2.B);
students missing item 2 need scaffolded ratio-formula practice (LO 2.2.A); students missing
item 3 need explicit connection to EK 2.2.B.3.
\end{skillbox}

\begin{skillbox}[Reinforcement \& Extension]{goldbox}
Homework (6 problems): equal-length interval test on two tables; construct a linear function
from two points; construct an exponential function from two points; contextual comparison of
linear vs.\ exponential growth over 12 months; AP-style MC identifying an exponential table;
extension proof that $b = (y_2/y_1)^{1/(x_2-x_1)}$ for any two points on $f(x)=ab^x$.

Preview of Lesson 2.3: the number $e$ and continuous exponential growth --- why the base
$e \approx 2.718$ arises naturally when compounding happens infinitely often.
\end{skillbox}

\end{document}
